\documentclass[14pt,a4paper]{extarticle}

\usepackage{fontspec}
\usepackage[russian]{babel}
\setmainfont{Times New Roman}
\setsansfont{Arial}
\setmonofont{Times New Roman}[Scale=1]

\usepackage[left=30mm,right=15mm,top=20mm,bottom=20mm]{geometry}
\usepackage{setspace}
\usepackage{indentfirst}
\usepackage{amsmath}
\usepackage{booktabs}
\usepackage{hhline}
\usepackage{longtable}
\usepackage{array}
\usepackage{tabularx}
\usepackage{enumitem}
\usepackage{listings}
\usepackage{xcolor}
\usepackage{tikz}
\usepackage{pgfplots}
\pgfplotsset{compat=1.18}
\usepackage{float}
\usepackage{caption}
\usepackage{pdfpages}
\usepackage{pdflscape}
\usepackage{hyperref}
\usepackage{titlesec}
\urlstyle{same}

\addto\captionsrussian{\renewcommand{\contentsname}{СОДЕРЖАНИЕ}}

\usetikzlibrary{arrows.meta,positioning,shapes.geometric}

\onehalfspacing
\setlength{\parindent}{1.25cm}
\setlength{\parskip}{0pt}
\setlength{\emergencystretch}{3em}
\hyphenpenalty=10000
\exhyphenpenalty=10000
\sloppy
\clubpenalty=10000
\widowpenalty=10000
\displaywidowpenalty=10000
\setlength{\rightskip}{0.5mm}
\AtBeginDocument{%
    \setlength{\abovedisplayskip}{0pt}%
    \setlength{\abovedisplayshortskip}{0pt}%
    \setlength{\belowdisplayskip}{0pt}%
    \setlength{\belowdisplayshortskip}{0pt}%
}
\hypersetup{
    hidelinks,
    pdftitle={Сравнительный анализ ресурсной эффективности алгоритмов},
    pdfauthor={Ананьев Никита Андреевич}
}

\setlist[itemize]{
    label=---,
    labelindent=\parindent,
    leftmargin=*,
    nosep
}
\setlist[enumerate]{
    label=\arabic*),
    labelindent=\parindent,
    leftmargin=*,
    nosep
}

\titleformat{\section}[hang]
    {\normalfont\Large\bfseries\raggedright\hyphenpenalty=10000\exhyphenpenalty=10000}
    {\hspace*{1.25cm}\thesection}{\fontdimen2\font}{}
\titleformat{name=\section,numberless}[block]
    {\normalfont\Large\bfseries\centering\hyphenpenalty=10000\exhyphenpenalty=10000}
    {}{0pt}{}
\titleformat{\subsection}[hang]
    {\normalfont\large\bfseries\raggedright\hyphenpenalty=10000\exhyphenpenalty=10000}
    {\hspace*{1.25cm}\thesubsection}{\fontdimen2\font}{}
\titlespacing*{\section}{0pt}
    {0pt}{2.3ex plus .2ex}
\titlespacing*{\subsection}{0pt}
    {3.25ex plus 1ex minus .2ex}{1.5ex plus .2ex}

\newcommand{\structuralpart}[1]{%
    \clearpage
    \section*{#1}%
    \phantomsection
    \addcontentsline{toc}{section}{#1}%
}

\captionsetup{labelsep=endash}
\captionsetup[figure]{name=Рисунок,justification=centering}
\captionsetup[table]{name=Таблица,justification=raggedright,
    singlelinecheck=false,skip=2pt}
\setlength{\doublerulesep}{0.5pt}

\makeatletter
\renewcommand{\tagform@}[1]{%
    \maketag@@@{\kern1mm(\ignorespaces#1\unskip\@@italiccorr)\kern1mm}}
% ссылка на формулу в тексте без отступов номера формулы
\renewcommand{\eqref}[1]{\textup{(\ref{#1})}}
\renewcommand{\@biblabel}[1]{#1.}
\renewenvironment{thebibliography}[1]
    {\section*{\refname}%
     \@mkboth{\MakeUppercase\refname}{\MakeUppercase\refname}%
     \list{\@biblabel{\@arabic\c@enumiv}}%
          {\settowidth\labelwidth{\@biblabel{#1}}%
           \leftmargin\z@
           \itemindent\dimexpr\parindent+\labelwidth+\labelsep\relax
           \rightmargin 5mm
           \@openbib@code
           \usecounter{enumiv}%
           \let\p@enumiv\@empty
           \renewcommand\theenumiv{\@arabic\c@enumiv}}%
     \sloppy
     \clubpenalty4000
     \@clubpenalty\clubpenalty
     \widowpenalty4000
     \sfcode`\.\@m}
    {\def\@noitemerr
         {\@latex@warning{Empty `thebibliography' environment}}%
     \endlist}
\renewcommand*\l@section[2]{%
    \ifnum\c@tocdepth>\z@
        \addpenalty\@secpenalty
        \addvspace{.5em \@plus\p@}%
        \setlength\@tempdima{\dimexpr\parindent+1.5em\relax}%
        \begingroup
        \parindent\z@ \rightskip\@pnumwidth
        \parfillskip-\@pnumwidth
        \leavevmode\bfseries
        \advance\leftskip\@tempdima
        \hskip-\leftskip
        \hspace*{1.25cm}%
        #1\nobreak\hfil\nobreak
        \hb@xt@\@pnumwidth{\hss #2}\par
        \endgroup
    \fi}
\renewcommand*\l@subsection[2]{%
    \@dottedtocline{2}{1.25cm}{2.8em}{#1}{#2}}
\makeatother

\lstset{
    basicstyle=\normalfont\normalsize\rmfamily\setstretch{1},
    keywordstyle=\color{black},
    commentstyle=\color{black},
    stringstyle=\color{black},
    numbers=left,
    numberstyle=\normalfont\normalsize\rmfamily,
    stepnumber=1,
    numbersep=5pt,
    xleftmargin=9mm,
    linewidth=\dimexpr\linewidth-9mm\relax,
    frame=single,
    breaklines=true,
    columns=fullflexible,
    keepspaces=true,
    showstringspaces=false,
    tabsize=4
}
\renewcommand{\lstlistingname}{Листинг}

\tikzset{
    algorithm block/.style={draw,align=center,minimum width=64mm,
        inner xsep=0pt,inner ysep=2pt},
    terminator/.style={algorithm block,rectangle,rounded corners=5mm,
        minimum height=10mm},
    data/.style={trapezium,trapezium left angle=70,
        trapezium right angle=110,trapezium stretches=true,
        algorithm block,text width=54mm,minimum height=10mm},
    process/.style={algorithm block,rectangle,text width=54mm,
        minimum height=10mm},
    decision/.style={diamond,draw,align=center,aspect=3.2,
        minimum width=64mm,minimum height=20mm,inner sep=0pt},
    loop start/.style={rectangle,align=center,draw=none,
        minimum width=64mm,minimum height=12mm,text width=54mm,
        inner xsep=0pt,inner ysep=2pt,path picture={
            \draw[black]
                ([yshift=-5mm]path picture bounding box.north west) --
                ([xshift=5mm]path picture bounding box.north west) --
                ([xshift=-5mm]path picture bounding box.north east) --
                ([yshift=-5mm]path picture bounding box.north east) --
                (path picture bounding box.south east) --
                (path picture bounding box.south west) -- cycle;
        }},
    loop end/.style={rectangle,align=center,draw=none,
        minimum width=64mm,minimum height=12mm,text width=54mm,
        inner xsep=0pt,inner ysep=2pt,path picture={
            \draw[black]
                (path picture bounding box.north west) --
                (path picture bounding box.north east) --
                ([yshift=5mm]path picture bounding box.south east) --
                ([xshift=-5mm]path picture bounding box.south east) --
                ([xshift=5mm]path picture bounding box.south west) --
                ([yshift=5mm]path picture bounding box.south west) -- cycle;
        }},
    flow/.style={draw,-{Latex[length=2.5mm]}},
    every node/.style={font=\small}
}


% Настройки ЛР № 3 (не входят в общую преамбулу)
\usepackage{multirow}
% Графовая модель или схема: рисунок TikZ, уменьшается, только если он
% шире текстового поля или выше 0,85 высоты страницы.
% строки таблицы из файла без лишней пустой строки в конце
\makeatletter\let\inputrows\@@input\makeatother
\newsavebox{\graphbox}
\newcommand{\fitgraph}[1]{%
    \sbox{\graphbox}{#1}%
    \ifdim\wd\graphbox>\linewidth
        \sbox{\graphbox}{\resizebox{\linewidth}{!}{\usebox{\graphbox}}}%
    \fi
    \ifdim\ht\graphbox>0.85\textheight
        \sbox{\graphbox}{\resizebox{!}{0.85\textheight}{\usebox{\graphbox}}}%
    \fi
    \usebox{\graphbox}}
\newcommand{\graphmodel}[1]{\fitgraph{\input{graphs/#1.tex}}}
\hypersetup{
    pdftitle={Рекурсивные и нерекурсивные алгоритмы},
    pdfauthor={Дудырев Д.~С.}
}
\graphicspath{{image/}}
\tikzset{
    gv/.style={circle,draw,minimum size=9mm,inner sep=0pt,font=\normalsize},
    hist/.style={draw,rounded corners=1.5mm,align=center,
        minimum width=22mm,minimum height=10mm,inner sep=1pt,font=\small},
    ret/.style={flow,dashed},
    lbl/.style={font=\small,inner sep=1pt},
    histnarrow/.style={hist,minimum width=17mm},
    junction/.style={circle,fill,inner sep=0pt,minimum size=1.5mm},
    % предопределённый процесс: дополнительные вертикальные линии у боковых сторон
    predefined/.style={process,path picture={
        \draw ([xshift=3mm]path picture bounding box.north west) --
              ([xshift=3mm]path picture bounding box.south west);
        \draw ([xshift=-3mm]path picture bounding box.north east) --
              ([xshift=-3mm]path picture bounding box.south east);}},
    % выравнивание по центру без растяжения пробелов внутри строки
    algorithm block/.append style={align=flush center},
    decision/.append style={align=flush center},
    % граница цикла: в v2 контур рисуется в path picture и обрезается по
    % самому контуру (видна половина толщины линии, стороны пропадают при
    % просмотре); здесь контур рисуется поверх узла без обрезки
    loop start/.style={rectangle,align=flush center,draw=none,
        minimum width=64mm,minimum height=12mm,text width=54mm,
        inner xsep=0pt,inner ysep=2pt,
        append after command={\pgfextra{\draw[black]
            ([yshift=-5mm]\tikzlastnode.north west) --
            ([xshift=5mm]\tikzlastnode.north west) --
            ([xshift=-5mm]\tikzlastnode.north east) --
            ([yshift=-5mm]\tikzlastnode.north east) --
            (\tikzlastnode.south east) --
            (\tikzlastnode.south west) -- cycle;}}},
    loop end/.style={rectangle,align=flush center,draw=none,
        minimum width=64mm,minimum height=12mm,text width=54mm,
        inner xsep=0pt,inner ysep=2pt,
        append after command={\pgfextra{\draw[black]
            (\tikzlastnode.north west) --
            (\tikzlastnode.north east) --
            ([yshift=5mm]\tikzlastnode.south east) --
            ([xshift=-5mm]\tikzlastnode.south east) --
            ([xshift=5mm]\tikzlastnode.south west) --
            ([yshift=5mm]\tikzlastnode.south west) -- cycle;}}},
    sprocess/.style={process,text width=26mm,minimum width=34mm},
    spredefined/.style={predefined,text width=26mm,minimum width=34mm},
    sdecision/.style={decision,minimum width=28mm,minimum height=13mm}
}
% Подпись листинга снизу; листинг не разрывается (окружение minipage)
\lstset{captionpos=b}
% Кириллица в листингах: без literate XeTeX выводит её вне очереди
% Переносы строк листинга выполняются через \discretionary, а преамбула v2
% запрещает их штрафом 10000; внутри листингов штраф снижается
% Код листингов — моноширинный Menlo, как в принятой ЛР № 1 (Times v2 в
% листингах сбивает отступы и заменяет * на \ast); Scale 0,86 даёт 12 пт
% (ГОСТ 7.32-2017, п. 6.1.1: не менее 12 пт). \texttt в тексте — Times (v2).
\newfontfamily\listingfont{Menlo}[Scale=0.86]
\lstset{basicstyle={\listingfont\normalsize\setstretch{1}\hyphenpenalty=50\relax\exhyphenpenalty=50\relax},
    columns=fixed,basewidth=0.6em,breakatwhitespace=true}
\lstset{extendedchars=true,literate={-}{{\char"2D}}1 {*}{{\char"2A}}1 {"}{{\textquotedbl}}1 {а}{{а}}1 {б}{{б}}1 {в}{{в}}1 {г}{{г}}1 {д}{{д}}1 {е}{{е}}1 {ё}{{ё}}1 {ж}{{ж}}1 {з}{{з}}1 {и}{{и}}1 {й}{{й}}1 {к}{{к}}1 {л}{{л}}1 {м}{{м}}1 {н}{{н}}1 {о}{{о}}1 {п}{{п}}1 {р}{{р}}1 {с}{{с}}1 {т}{{т}}1 {у}{{у}}1 {ф}{{ф}}1 {х}{{х}}1 {ц}{{ц}}1 {ч}{{ч}}1 {ш}{{ш}}1 {щ}{{щ}}1 {ъ}{{ъ}}1 {ы}{{ы}}1 {ь}{{ь}}1 {э}{{э}}1 {ю}{{ю}}1 {я}{{я}}1 {А}{{А}}1 {Б}{{Б}}1 {В}{{В}}1 {Г}{{Г}}1 {Д}{{Д}}1 {Е}{{Е}}1 {Ё}{{Ё}}1 {Ж}{{Ж}}1 {З}{{З}}1 {И}{{И}}1 {Й}{{Й}}1 {К}{{К}}1 {Л}{{Л}}1 {М}{{М}}1 {Н}{{Н}}1 {О}{{О}}1 {П}{{П}}1 {Р}{{Р}}1 {С}{{С}}1 {Т}{{Т}}1 {У}{{У}}1 {Ф}{{Ф}}1 {Х}{{Х}}1 {Ц}{{Ц}}1 {Ч}{{Ч}}1 {Ш}{{Ш}}1 {Щ}{{Щ}}1 {Ъ}{{Ъ}}1 {Ы}{{Ы}}1 {Ь}{{Ь}}1 {Э}{{Э}}1 {Ю}{{Ю}}1 {Я}{{Я}}1 {–}{{–}}1 {—}{{—}}1 {«}{{«}}1 {»}{{»}}1 {№}{{№}}1}
\lstdefinelanguage{Swift}{
    morekeywords={func,let,var,if,else,for,in,return,import,public,struct,
        break,guard,static,true,false,nil},
    sensitive=true,
    morecomment=[l]{//},
    morestring=[b]"
}

\begin{document}
\begin{titlepage}
\noindent\begin{minipage}{0.15\textwidth}
    \includegraphics[width=\linewidth]{bmstu-logo.pdf}
\end{minipage}%
\begin{minipage}{0.84\textwidth}\centering\small\singlespacing
    \textbf{Министерство науки и высшего образования Российской Федерации}\\
    \textbf{Федеральное государственное автономное образовательное учреждение высшего образования}\\
    \textbf{«Московский государственный технический университет имени Н.~Э.~Баумана}\\
    \textbf{(национальный исследовательский университет)»}\\
    \textbf{(МГТУ им. Н.~Э.~Баумана)}
\end{minipage}

{\rightskip=0pt\noindent\rule{\linewidth}{3pt}\par}

\vspace{5mm}
% ровно 14 пт (основной текст 14,4 пт), иначе строка кафедры не помещается в поля
{\fontsize{14pt}{17pt}\selectfont
\noindent ФАКУЛЬТЕТ \underline{«Информатика и системы управления»}\par
\vspace{3mm}
\noindent КАФЕДРА \underline{«Программное обеспечение ЭВМ и информационные технологии»}\par}

\vspace{2cm}
\begin{center}
    {\Large\textbf{Лабораторная работа № 3}\par}
    \vspace{2mm}
    {\Large\textbf{по дисциплине «Анализ алгоритмов»}\par}
\end{center}

\vspace{2cm}
\noindent\textbf{Тема} \underline{Рекурсивные и нерекурсивные алгоритмы}\par\vspace{4mm}
\noindent\textbf{Студент} \underline{Дудырев Д.~С.}\par\vspace{4mm}
\noindent\textbf{Группа} \underline{ИУ7-52Б}\par\vspace{4mm}
\noindent\textbf{Преподаватель} \underline{Волкова Л.~Л.}\par\vspace{4mm}
\noindent\textbf{Преподаватель} \underline{Строганов Д.~М.}\par

\vfill
\begin{center}
    Москва, \the\year
\end{center}
\end{titlepage}
\setcounter{page}{2}

\tableofcontents

\structuralpart{ВВЕДЕНИЕ}

Рекурсивным называется алгоритм, который в процессе выполнения обращается
к самому себе непосредственно или через другие алгоритмы. Нерекурсивный
алгоритм обходится без таких обращений и повторяет действия с помощью циклов.
Графовой моделью алгоритма называется ориентированный граф, вершины которого
соответствуют операторам алгоритма или их однократным выполнениям,
а дуги~--- связям между ними по передаче управления или данных.

Цель работы~--- провести сравнение рекурсивного и нерекурсивного алгоритмов
на основе графовых моделей алгоритмов.

Для достижения цели необходимо решить следующие задачи:
\begin{itemize}
    \item дать определения четырёх графовых моделей алгоритмов,
    операционного и информационного отношений;
    \item разработать нерекурсивный и рекурсивный последовательные алгоритмы
    решения задачи варианта;
    \item реализовать разработанные алгоритмы;
    \item провести тестирование реализаций алгоритмов;
    \item описать каждый алгоритм четырьмя графовыми моделями;
    \item обосновать, какие фрагменты алгоритмов или их реализаций могут быть
    выполнены в параллельном режиме, а какие не могут.
\end{itemize}

Задача варианта~--- определение количества $k$-значных натуральных чисел
с заданной суммой цифр $s$.

\clearpage
\section{Аналитическая часть}

\subsection{Постановка задачи}

Даны натуральные числа $k$ и $s$. Требуется определить, сколько существует
$k$-значных натуральных чисел, сумма цифр которых равна $s$. Запись натурального
числа не может начинаться с цифры 0. Например, при $k=3$ и $s=15$ ответ равен 69.
Если $s>9k$, искомых чисел нет, поскольку каждая цифра не превосходит 9.

Ответ не превосходит количества всех $k$-значных чисел $9\cdot10^{k-1}$.
При $k\le19$ это значение меньше $2^{63}$, поэтому ответ представим
64-битовым целым числом со знаком; для б\'{о}льших $k$ такое представление
может переполниться.

Нерекурсивный алгоритм решает задачу методом динамического
программирования: ответ составляется
из ответов на подзадачи меньшего размера, каждая подзадача решается один раз,
а её результат сохраняется для повторного использования~\cite{laaksonen}.
Обозначим через $F(i,j)$ количество $i$-значных натуральных чисел, сумма цифр
которых равна~$j$. Каждое $i$-значное число при $i\ge2$ получается приписыванием
справа цифры $c$ к $(i-1)$-значному числу с суммой цифр $j-c$, откуда
\begin{equation}
    \label{eq:f}
    F(i,j)=
    \begin{cases}
        1, & i=0,\ j=0,\\
        0, & i=0,\ j>0,\\
        0, & i\ge1,\ j=0,\\
        \sum\limits_{c=0}^{\min(9,j)}F(i-1,j-c), & i\ge1,\ j\ge1,
    \end{cases}
\end{equation}
где $c$~--- последняя цифра числа. Значение $F(0,0)=1$ соответствует пустой
записи. При $i=1$ в сумме отлична от нуля только величина $F(0,0)$, поэтому
$F(1,j)=1$ для $1\le j\le9$, а первая цифра числа всегда отлична от нуля.
Ответ задачи~--- $F(k,s)$.

На рекурсивный алгоритм наложены ограничения задания: нельзя использовать циклы,
массивы, строки и срезы. Поэтому сумма в формуле~\eqref{eq:f} заменяется
рекурсивной функцией
\begin{equation}
    \label{eq:s}
    S(i,j,c)=
    \begin{cases}
        0, & c>9 \text{ или } c>j,\\
        F(i-1,j-c)+S(i,j,c+1), & \text{иначе},
    \end{cases}
\end{equation}
где $S(i,j,c)$~--- сумма величин $F(i-1,j-d)$ для цифр $d$ от $c$ до $\min(9,j)$.
Тогда при $i\ge1$ и $j\ge1$ выполняется $F(i,j)=S(i,j,0)$.

\subsection{Графовые модели алгоритмов}

Алгоритм рассматривается как совокупность операторов. Однократное выполнение
оператора называется его срабатыванием. Между операторами алгоритма определяются
два отношения~\cite{voevodin}:
\begin{itemize}
    \item операционное отношение: оператор $v$ связан с оператором $u$, если $v$
    может быть выполнен непосредственно после $u$, то есть управление
    передаётся от $u$ к $v$;
    \item информационное отношение: оператор $v$ связан с оператором $u$, если $v$
    использует в качестве аргумента величину, вычисленную $u$, и эта величина
    не изменяется между их срабатываниями.
\end{itemize}

На основе этих отношений строятся четыре графовые модели~\cite{voevodin}:
\begin{itemize}
    \item граф управления~--- вершины соответствуют операторам, дуги~---
    операционному отношению;
    \item информационный граф~--- вершины соответствуют операторам, дуги~---
    информационному отношению;
    \item операционная история~--- вершины соответствуют срабатываниям операторов
    при конкретных входных данных, дуги указывают порядок срабатываний;
    \item информационная история~--- вершины соответствуют срабатываниям
    операторов при конкретных входных данных, дуги~--- информационному отношению
    между срабатываниями.
\end{itemize}

Граф управления и информационный граф описывают все возможные связи
операторов, а истории~--- одно выполнение алгоритма. Информационная история
не содержит циклов. Ярусом её вершины называется число вершин на наибольшем
пути, оканчивающемся в этой вершине. Срабатывания одного яруса не связаны
информационным отношением и могут выполняться параллельно, а число ярусов
равно наименьшему числу последовательных шагов при неограниченном числе
исполнителей и одинаковой длительности срабатываний.

\clearpage
\section{Конструкторская часть}

\subsection{Требования к программе}

К программе предъявляются следующие требования:
\begin{itemize}
    \item реализации алгоритмов принимают натуральные числа $s$ и $k$, \mbox{$k\le19$},
    и возвращают количество $k$-значных натуральных чисел с суммой цифр $s$;
    \item рекурсивная реализация не использует циклы, массивы, строки и срезы;
    \item программа выводит таблицу значений $F(i,j)$ для \mbox{$0\le i\le k$},
    \mbox{$0\le j\le s$} и результаты обоих алгоритмов.
\end{itemize}

\subsection{Нерекурсивный алгоритм}

Нерекурсивный алгоритм заполняет таблицу значений $F$ по формуле~\eqref{eq:f}:
строки~--- сверху вниз по возрастанию $i$, ячейки строки~--- слева направо
по возрастанию $j$. Алгоритм получает числа $k$ и $s$ и возвращает $F[k][s]$.
Псевдокод алгоритма приведён в листинге~\ref{lst:iter} «Псевдокод
нерекурсивного алгоритма», схема показана на рисунке~\ref{fig:scheme-iter}.

\bigskip
\noindent\begin{minipage}{\dimexpr\linewidth-0.5mm\relax}
\begin{lstlisting}[language=Python,caption={Псевдокод нерекурсивного алгоритма},label={lst:iter}]
def count(k, s):
    F = [[0] * (s + 1) for _ in range(k + 1)]
    F[0][0] = 1
    for i in range(1, k + 1):
        for j in range(1, s + 1):
            for c in range(0, min(9, j) + 1):
                F[i][j] += F[i - 1][j - c]
    return F[k][s]
\end{lstlisting}
\end{minipage}
\medskip

\begin{figure}[H]
    \centering
    \begin{tikzpicture}[node distance=5mm]
\node[terminator] (start) {Начало};
\node[process,below=of start] (b2) {Создать таблицу $F$ размером $(k+1)\times(s+1)$, заполненную нулями};
\node[process,below=of b2] (b3) {$F[0][0]=1$};
\node[loop start,below=of b3] (b4) {Цикл I\\$i$ от 1 до $k$};
\node[loop start,below=of b4] (b5) {Цикл J\\$j$ от 1 до $s$};
\node[loop start,below=of b5] (b6) {Цикл C\\$c$ от 0 до $\min(9,j)$};
\node[process,below=of b6] (b7) {$F[i][j]=F[i][j]+F[i-1][j-c]$};
\node[loop end,below=of b7] (e6) {Цикл C};
\node[loop end,below=of e6] (e5) {Цикл J};
\node[loop end,below=of e5] (e4) {Цикл I};
\node[process,below=of e4] (b8) {Вернуть $F[k][s]$};
\node[terminator,below=of b8] (end) {Конец};
\foreach \a/\b in {start/b2,b2/b3,b3/b4,b4/b5,b5/b6,b6/b7,b7/e6,e6/e5,e5/e4,e4/b8,b8/end}
    \draw[flow] (\a) -- (\b);
\end{tikzpicture}

    \caption{Схема нерекурсивного алгоритма}
    \label{fig:scheme-iter}
\end{figure}

\subsection{Рекурсивный алгоритм}

Рекурсивный алгоритм вычисляет $F(k,s)$ по формулам~\eqref{eq:f}
и~\eqref{eq:s} без сохранения промежуточных значений: перебор последней цифры
выполняется рекурсивной функцией $S$. Схемы функций $F$ и $S$ показаны
на рисунках~\ref{fig:scheme-f} и~\ref{fig:scheme-s}, псевдокод обеих функций
приведён в листинге~\ref{lst:rec} «Псевдокод рекурсивного алгоритма».

\begin{figure}[H]
    \centering
    \begin{tikzpicture}[x=1mm,y=-1mm]
\node[terminator] (start) at (0,0) {$F(i,j)$};
\node[sdecision] (b1) at (0,22) {$i=0$};
\node[sdecision] (b2) at (-40,44) {$j=0$};
\node[sdecision] (b5) at (40,44) {$j=0$};
\node[sprocess] (b3) at (-60,66) {$r=1$};
\node[sprocess] (b4) at (-20,66) {$r=0$};
\node[sprocess] (b6) at (20,66) {$r=0$};
\node[spredefined] (b7) at (60,66) {$r=S(i,j,0)$};
\node[terminator] (end) at (0,104) {Возврат: $r$};
\draw[flow] (start) -- (b1);
\draw[flow] (b1.west) -- node[lbl,above] {да} (-40,22) -- (b2.north);
\draw[flow] (b1.east) -- node[lbl,above] {нет} (40,22) -- (b5.north);
\draw[flow] (b2.west) -- (-60,44) -- node[lbl,left,pos=0.35] {да} (b3.north);
\draw[flow] (b2.east) -- (-20,44) -- node[lbl,right,pos=0.35] {нет} (b4.north);
\draw[flow] (b5.west) -- (20,44) -- node[lbl,left,pos=0.35] {да} (b6.north);
\draw[flow] (b5.east) -- (60,44) -- node[lbl,right,pos=0.35] {нет} (b7.north);
% слияние линий в смещённых точках (ГОСТ 19.701-90, п. 4.2.3)
\draw (b3.south) |- (0,77);
\draw[flow] (b4.south) |- (0,81);
\draw[flow] (b6.south) |- (0,85);
\draw[flow] (b7.south) |- (0,89);
\draw[flow] (0,77) -- (end.north);
\end{tikzpicture}

    \caption{Схема алгоритма вычисления функции $F$}
    \label{fig:scheme-f}
\end{figure}

\begin{figure}[H]
    \centering
    \begin{tikzpicture}[x=1mm,y=-1mm]
\node[terminator] (start) at (0,0) {$S(i,j,c)$};
\node[decision] (b9) at (0,24) {$c>9$ или $c>j$};
\node[process] (b10) at (-38,50) {$t=0$};
\node[predefined] (b11) at (38,50) {$a=F(i-1,j-c)$};
\node[predefined] (b12) at (38,68) {$b=S(i,j,c+1)$};
\node[process] (b13) at (38,86) {$t=a+b$};
\node[terminator] (end) at (0,116) {Возврат: $t$};
\draw[flow] (start) -- (b9);
\draw[flow] (b9.west) -- node[lbl,above] {да} (-38,24) -- (b10.north);
\draw[flow] (b9.east) -- node[lbl,above] {нет} (38,24) -- (b11.north);
\draw[flow] (b11) -- (b12);
\draw[flow] (b12) -- (b13);
% слияние линий в смещённых точках (ГОСТ 19.701-90, п. 4.2.3)
\draw (b10.south) |- (0,97);
\draw[flow] (b13.south) |- (0,102);
\draw[flow] (0,97) -- (end.north);
\end{tikzpicture}

    \caption{Схема алгоритма вычисления функции $S$}
    \label{fig:scheme-s}
\end{figure}

\bigskip
\noindent\begin{minipage}{\dimexpr\linewidth-0.5mm\relax}
\begin{lstlisting}[language=Python,caption={Псевдокод рекурсивного алгоритма},label={lst:rec}]
def F(i, j):
    if i == 0:
        if j == 0:
            r = 1
        else:
            r = 0
    elif j == 0:
        r = 0
    else:
        r = S(i, j, 0)
    return r
def S(i, j, c):
    if c > 9 or c > j:
        t = 0
    else:
        a = F(i - 1, j - c)
        b = S(i, j, c + 1)
        t = a + b
    return t
\end{lstlisting}
\end{minipage}
\medskip

Вызов функции обозначен на схемах символом «Предопределённый процесс»
(строки 10, 16 и 17 листинга~\ref{lst:rec}). Выполнение вызова передаёт
аргументы вызываемой функции, а результат, возвращённый вызванной функцией
в строке 11 или 19, присваивается переменной $r$, $a$ или $b$. Повторные вызовы $F$
с одинаковыми аргументами вычисляются заново, поскольку сохранять результаты
в массиве запрещено.

В графовых моделях исследовательской части оператор обозначается номером
строки листинга~\ref{lst:iter} или~\ref{lst:rec}. Строки с заголовками
функций и строки \texttt{else} операторами не считаются.

\clearpage
\section{Технологическая часть}

\subsection{Средства реализации}

Алгоритмы реализованы на языке Свифт (Swift) версии 6.4.2. При выполнении
программы язык проверяет переполнение целочисленных операций и аварийно
завершает программу вместо получения неверного результата~\cite{swift}.
Для сборки используется диспетчер пакетов Свифт (Swift Package Manager),
для тестирования~--- библиотека Свифт Тестинг (Swift Testing)~\cite{swifttesting}.
Программа собиралась и выполнялась под управлением операционной системы
макОС (macOS) 26.6.2 на компьютере с процессором Apple M1 Pro и оперативной
памятью 16~ГБ.

\subsection{Реализация алгоритмов}

Реализации алгоритмов приведены в приложении~А. Нерекурсивный алгоритм
реализован функцией \texttt{countNumbersBySumDigits}
(листинг~\ref{lst:impl-iter} «Реализация нерекурсивного алгоритма»).
В ней цифры перебираются от 0 до 9 с досрочным выходом при отрицательной
сумме $j-c$, что равносильно границе $\min(9,j)$ в листинге~\ref{lst:iter}.

Функция \texttt{countNumbersBySumDigitsRecursive} реализует рекурсивный
алгоритм (листинг~\ref{lst:impl-rec} «Реализация рекурсивного алгоритма»).
Она вызывает вспомогательные функции,
соответствующие функциям $F$ и $S$ листинга~\ref{lst:rec}.

Для вывода таблицы значений программа использует функцию
\texttt{countNumbersBySumDigitsMatrix}, которая возвращает таблицу целиком
(листинг~\ref{lst:impl-table} «Формирование и вывод таблицы значений»).

\subsection{Сборка и запуск}

Команды сборки, запуска программы и тестирования выполняются из каталога
\texttt{lab-03} и приведены в листинге~\ref{lst:run} «Команды сборки, запуска
и тестирования».

\bigskip
\noindent\begin{minipage}{\dimexpr\linewidth-0.5mm\relax}
\begin{lstlisting}[caption={Команды сборки, запуска и тестирования},label={lst:run}]
swift build
swift run lab03
swift test
\end{lstlisting}
\end{minipage}
\medskip

При запуске программа выводит таблицу значений $F(i,j)$ для $s=15$, $k=5$
и результаты обоих алгоритмов: 2654.

\subsection{Тестирование}

Функциональные тесты приведены в листинге~\ref{lst:tests}
«Функциональные тесты» приложения~А. Одинаковые входные данные и ожидаемые
результаты для обеих реализаций приведены в таблице~\ref{tbl:tests}.

\begin{table}[H]
\caption{Функциональные тесты}
\label{tbl:tests}
\setlength{\extrarowheight}{2pt}
\begin{tabularx}{\linewidth}{|c|c|c|X|}
\hline
$s$ & $k$ & Ожидаемый результат & Проверяемый случай \\
\hhline{|=:=:=:=|}
4 & 1 & 1 & однозначное число \\ \hline
14 & 1 & 0 & сумма больше $9k$ для однозначных чисел \\ \hline
15 & 3 & 69 & пример из условия задачи \\ \hline
28 & 3 & 0 & сумма больше $9k$ \\ \hline
2 & 2 & 2 & числа 11 и 20 \\ \hline
1 & 4 & 1 & наименьшая сумма, число 1000 \\ \hline
36 & 4 & 1 & наибольшая сумма, число 9999 \\ \hline
\end{tabularx}
\end{table}

Все 14 проверок (по 7 для каждой реализации) пройдены успешно.

\clearpage
\section{Исследовательская часть}

Графовые модели построены по псевдокодам листингов~\ref{lst:iter}
и~\ref{lst:rec}. Вершина графа управления и информационного графа обозначена
номером строки оператора. Вершина истории содержит номер строки оператора
и пояснение срабатывания. Истории построены для входных данных $k=1$, $s=1$,
при которых ответ равен 1 (число 1).

\subsection{Графовые модели нерекурсивного алгоритма}

Граф управления нерекурсивного алгоритма показан на рисунке~\ref{fig:iter-control}.
Из начала цикла управление передаётся либо в тело цикла, либо
к оператору, следующему за циклом.

\begin{figure}[H]
    \centering
    \graphmodel{iter-control}
    \caption{Граф управления нерекурсивного алгоритма}
    \label{fig:iter-control}
\end{figure}

Информационный граф нерекурсивного алгоритма показан на
рисунке~\ref{fig:iter-info}. Петли у вершин 4--6 соответствуют вычислению
следующего значения переменной цикла по предыдущему, петля у вершины 7~---
накоплению суммы в ячейке таблицы.

\begin{figure}[H]
    \centering
    \graphmodel{iter-info}
    \caption{Информационный граф нерекурсивного алгоритма}
    \label{fig:iter-info}
\end{figure}

Операционная история нерекурсивного алгоритма показана на
рисунке~\ref{fig:iter-op-history}. Для начала цикла указано значение
переменной цикла; срабатывание, при котором условие продолжения ложно, является
последним срабатыванием начала цикла. Для оператора 7 указана изменяемая
ячейка таблицы.

\begin{figure}[H]
    \centering
    \graphmodel{iter-op-history}
    \caption{Операционная история нерекурсивного алгоритма}
    \label{fig:iter-op-history}
\end{figure}

Информационная история нерекурсивного алгоритма показана на
рисунке~\ref{fig:iter-info-history}. Срабатывания операторов 2--5 в начале
выполнения не используют результатов других срабатываний.

\begin{figure}[H]
    \centering
    \graphmodel{iter-info-history}
    \caption{Информационная история нерекурсивного алгоритма}
    \label{fig:iter-info-history}
\end{figure}

\subsection{Графовые модели рекурсивного алгоритма}

Граф управления рекурсивного алгоритма показан на рисунке~\ref{fig:rec-control}.
Сплошные дуги от операторов 10, 16 и 17 соответствуют вызовам функций,
штриховые дуги от операторов 11 и 19~--- возврату управления из функции
к оператору вызывающей функции, который выполняется после вызова.

\begin{figure}[H]
    \centering
    \graphmodel{rec-control}
    \caption{Граф управления рекурсивного алгоритма}
    \label{fig:rec-control}
\end{figure}

Информационный граф рекурсивного алгоритма показан на
рисунке~\ref{fig:rec-info}. Дуги от операторов вызова 10, 16 и 17 соответствуют
передаче аргументов, дуги от операторов 11 и 19~--- передаче результата
функции. Точками отмечено разветвление общей линии дуг от оператора 16.

\begin{figure}[H]
    \centering
    \graphmodel{rec-info}
    \caption{Информационный граф рекурсивного алгоритма}
    \label{fig:rec-info}
\end{figure}

Операционная история рекурсивного алгоритма показана на
рисунке~\ref{fig:rec-op-history}. Для каждого срабатывания указан вызов
функции, в котором оно выполнено.

\begin{figure}[H]
    \centering
    \graphmodel{rec-op-history}
    \caption{Операционная история рекурсивного алгоритма}
    \label{fig:rec-op-history}
\end{figure}

Информационная история рекурсивного алгоритма показана на
рисунке~\ref{fig:rec-info-history}. Срабатывания операторов 2, 7 и 10 в вызове
$F(1,1)$ используют только входные данные, а срабатывания операторов 4, 6
и 14 присваивают постоянные значения, поэтому входящих дуг у них нет.

\begin{figure}[H]
    \centering
    \graphmodel{rec-info-history}
    \caption{Информационная история рекурсивного алгоритма}
    \label{fig:rec-info-history}
\end{figure}

\subsection{Сравнение алгоритмов по графовым моделям}

Число вершин операционной истории равно числу срабатываний операторов,
число ярусов информационной истории~--- наименьшему числу последовательных
шагов при неограниченном числе исполнителей. Эти характеристики для
нескольких входных данных приведены в таблице~\ref{tbl:stats}; они получены
построением историй по псевдокодам листингов~\ref{lst:iter} и~\ref{lst:rec}.

\begin{table}[H]
\caption{Характеристики историй алгоритмов}
\label{tbl:stats}
\setlength{\extrarowheight}{2pt}
\begin{tabularx}{\linewidth}{|c|c|>{\centering\arraybackslash}X|>{\centering\arraybackslash}X|>{\centering\arraybackslash}X|>{\centering\arraybackslash}X|}
\hline
\multirow{2}{*}{$k$} & \multirow{2}{*}{$s$} &
\multicolumn{2}{c|}{Нерекурсивный алгоритм} &
\multicolumn{2}{c|}{Рекурсивный алгоритм} \\
\cline{3-6}
& & срабатываний & ярусов & срабатываний & ярусов \\
\hhline{|=:=:=:=:=:=|}
\inputrows graphs/stats.tex \end{tabularx}
\end{table}

Число срабатываний нерекурсивного алгоритма пропорционально числу
слагаемых в формуле~\eqref{eq:f}, то есть не превосходит величины порядка
$10ks$. В рекурсивном алгоритме значения $F(i,j)$ с одинаковыми аргументами
вычисляются повторно, поэтому число срабатываний растёт значительно быстрее:
при $k=4$, $s=18$ оно почти в 45 раз больше, чем у нерекурсивного алгоритма.
Число ярусов информационной истории у алгоритмов близко, а при б\'{о}льших
входных данных у рекурсивного алгоритма оно меньше. У нерекурсивного
алгоритма его увеличивают цепочки срабатываний начала цикла, каждое
из которых использует предыдущее значение переменной цикла.

\subsection{Фрагменты, допускающие параллельное выполнение}

В нерекурсивном алгоритме параллельно могут выполняться:
\begin{itemize}
    \item заполнение нулями разных ячеек таблицы в операторе 2;
    \item вычисление ячеек одной строки таблицы при фиксированном $i$: значение
    $F(i,j)$ использует только строку $i-1$ и записывается в собственную ячейку,
    поэтому итерации цикла по $j$ информационно независимы;
    \item сложение до десяти слагаемых суммы в формуле~\eqref{eq:f}, если
    накапливать частичные суммы отдельно и складывать их попарно; при таком
    малом числе слагаемых выигрыш не покрывает затрат на организацию
    параллельного выполнения.
\end{itemize}

В нерекурсивном алгоритме не могут выполняться параллельно:
\begin{itemize}
    \item вычисление разных строк таблицы: строка $i$ использует значения
    строки $i-1$, что видно по дугам информационной истории от срабатываний
    операторов 2 и 3, записавших строку 0, к срабатываниям оператора 7
    для ячейки $F[1][1]$;
    \item прибавление слагаемых к одной ячейке в операторе 7 без отдельных
    частичных сумм, так как каждое прибавление использует результат
    предыдущего;
    \item оператор 3 выполняется после заполнения ячейки $F[0][0]$ нулями,
    иначе записанное им значение будет потеряно;
    \item оператор 8 выполняется после вычисления ячейки $F[k][s]$.
\end{itemize}

Цепочки срабатываний начала цикла не ограничивают параллельное
выполнение: в параллельной реализации индексы ячеек могут быть заданы
каждому исполнителю непосредственно.

В рекурсивном алгоритме параллельно могут выполняться вызовы 16
и 17 одного вызова $S$: оба используют только аргументы этого вызова
и не связаны информационным отношением, что видно по общему предшественнику
и одинаковому ярусу их срабатываний в информационной истории. Тем самым
независимыми являются все ветви дерева рекурсивных вызовов.

В рекурсивном алгоритме не могут выполняться параллельно:
\begin{itemize}
    \item оператор 18 и вызовы 16 и 17, поскольку сумма требует
    результатов обоих вызовов;
    \item оператор 11 и вызов 10, поскольку оператор 11 возвращает результат вызова;
    \item проверки 2, 3, 7 и 13 и следующие за ними операторы, так как
    от результата проверки зависит выполнение остальных операторов вызова.
\end{itemize}

Каждый вызов выполняет лишь несколько операторов, поэтому при параллельной
реализации создание исполнителя для каждого вызова требует больше времени,
чем сам вызов; параллельно целесообразно выполнять только верхние уровни
дерева вызовов.

\structuralpart{ЗАКЛЮЧЕНИЕ}

Цель работы достигнута: рекурсивный и нерекурсивный алгоритмы определения
количества $k$-значных натуральных чисел с заданной суммой цифр сравнены
на основе графовых моделей.

В ходе выполнения работы решены следующие задачи:
\begin{itemize}
    \item даны определения операционного и информационного отношений, графа
    управления, информационного графа, операционной и информационной историй;
    \item разработаны нерекурсивный алгоритм заполнения таблицы значений
    и рекурсивный алгоритм без циклов, массивов, строк и срезов;
    \item алгоритмы реализованы на языке Свифт (Swift);
    \item реализации прошли все 14 функциональных проверок, по 7 для каждой
    реализации;
    \item каждый алгоритм описан четырьмя графовыми моделями; истории построены
    для $k=1$, $s=1$;
    \item показано, что в нерекурсивном алгоритме параллельно могут вычисляться
    ячейки одной строки таблицы, а строки~--- только последовательно;
    в рекурсивном алгоритме независимы вызовы $F(i-1,j-c)$ и $S(i,j,c+1)$,
    а сумма их результатов выполняется после обоих вызовов.
\end{itemize}

По графовым моделям рекурсивный алгоритм требует значительно большего числа
срабатываний операторов из-за повторного вычисления одинаковых значений,
а число ярусов информационных историй обоих алгоритмов близко.

\clearpage
\phantomsection
\addcontentsline{toc}{section}{СПИСОК ИСПОЛЬЗОВАННЫХ ИСТОЧНИКОВ}
\renewcommand{\refname}{СПИСОК ИСПОЛЬЗОВАННЫХ ИСТОЧНИКОВ}
\begin{thebibliography}{9}
\bibitem{voevodin}
Воеводин,~В.~В. Параллельные вычисления / В.~В.~Воеводин, Вл.~В.~Воеводин.~---
Санкт-Петербург~: БХВ-Петербург, 2002.~--- 608~с.

\bibitem{laaksonen}
Laaksonen,~A. Guide to Competitive Programming~: Learning and Improving
Algorithms Through Contests / A.~Laaksonen.~--- 3rd~ed.~--- Cham~: Springer,
2024.~--- 349~p.~--- ISBN 978-3-031-61793-5.~--- DOI 10.1007/978-3-031-61794-2.

\bibitem{swift}
The Swift Programming Language~: документация~// Swift.org~: [сайт].~---
URL: \url{https://docs.swift.org/swift-book/} (дата обращения: 02.11.2026).

\bibitem{swifttesting}
Swift Testing~: документация~// Apple Developer~: [сайт].~---
URL: \url{https://developer.apple.com/documentation/testing}
(дата обращения: 02.11.2026).
\end{thebibliography}

\clearpage
\section*{ПРИЛОЖЕНИЕ А}
\phantomsection
\addcontentsline{toc}{section}{ПРИЛОЖЕНИЕ А}
\begin{center}
    \large\bfseries Исходный код реализаций
\end{center}

\setcounter{lstlisting}{0}
\renewcommand{\thelstlisting}{А.\arabic{lstlisting}}

\noindent\begin{minipage}{\dimexpr\linewidth-0.5mm\relax}
\lstinputlisting[language=Swift,caption={Реализация нерекурсивного алгоритма},label={lst:impl-iter}]{code/iterative.swift}
\end{minipage}

\noindent\begin{minipage}{\dimexpr\linewidth-0.5mm\relax}
\lstinputlisting[language=Swift,caption={Реализация рекурсивного алгоритма},label={lst:impl-rec}]{code/recursive.swift}
\end{minipage}

\noindent\begin{minipage}{\dimexpr\linewidth-0.5mm\relax}
\lstinputlisting[language=Swift,caption={Формирование и вывод таблицы значений},label={lst:impl-table}]{code/table.swift}
\end{minipage}

\noindent\begin{minipage}{\dimexpr\linewidth-0.5mm\relax}
\lstinputlisting[language=Swift,caption={Функциональные тесты},label={lst:tests}]{code/tests.swift}
\end{minipage}


\end{document}
